\chapter{Standard Library and System Calls}

\section{Introduction}

Every C programmer uses functions like \texttt{printf}, \texttt{malloc}, \texttt{open}, and \texttt{read}. But where do these functions come from? Are they part of the C language, part of the operating system, or something else entirely?

This chapter explores the layered architecture that sits between your C code and the bare hardware. We'll distinguish between:
\begin{itemize}
    \item \textbf{The C standard library (libc)}: User-space code that provides the standard C API
    \item \textbf{System calls}: The kernel interface that performs privileged operations
    \item \textbf{The boundary between user mode and kernel mode}: Where protection and privilege transitions occur
\end{itemize}

Understanding this distinction is crucial for:
\begin{itemize}
    \item \textbf{Performance optimization}: Knowing when library functions vs. raw syscalls are appropriate
    \item \textbf{Debugging}: Understanding error reporting and stack traces
    \item \textbf{Portability}: Writing code that works across different platforms
    \item \textbf{Security}: Understanding privilege boundaries and potential vulnerabilities
\end{itemize}

\section{The Layered Architecture}

Modern Unix-like systems have a clear separation between user-space programs and the kernel:

\begin{verbatim}
┌─────────────────────────────────────────────────────────────┐
│                    Your C Program                           │
│                                                             │
│  main() {                                                   │
│      printf("Hello\n");  // ← Not a syscall!                │
│      malloc(1024);        // ← Not a syscall!               │
│  }                                                          │
└─────────────────────────────────────────────────────────────┘
                              │
                              ▼
┌─────────────────────────────────────────────────────────────┐
│              C Standard Library (libc)                      │
│                                                             │
│  • Wrappers for standard C functions                        │
│  • Buffering, caching, optimization                         │
│  • POSIX API implementation                                 │
│  • Error handling (errno)                                   │
│                                                             │
│  Examples: printf(), malloc(), fopen(), strlen()            │
└─────────────────────────────────────────────────────────────┘
                              │
                              ▼
┌─────────────────────────────────────────────────────────────┐
│                 System Call Interface                       │
│                                                             │
│  • Well-defined entry points into kernel                    │
│  • Mode switch: user → kernel                               │
│  • Argument validation and copying                          │
│                                                             │
│  Examples: write(), brk(), mmap(), open()                   │
└─────────────────────────────────────────────────────────────┘
                              │
                              ▼
┌─────────────────────────────────────────────────────────────┐
│                      Operating System Kernel                │
│                                                             │
│  • Process management                                       │
│  • Memory management                                        │
│  • File system                                              │
│  • Device drivers                                           │
│  • Security and access control                              │
└─────────────────────────────────────────────────────────────┘
\end{verbatim}

\textbf{Key insight}: Most C library functions are \textbf{not} system calls. They're user-space code that may or may not eventually make system calls.

\section{What is a System Call?}

A \textbf{system call} (syscall) is a programmatic way a program requests a service from the operating system kernel. System calls are the only legitimate way to access kernel resources.

\subsection{Why System Calls Exist}

User programs cannot directly access:
\begin{itemize}
    \item \textbf{Hardware devices}: Direct hardware access would allow any program to monopolize or damage devices
    \item \textbf{Other processes' memory}: Without isolation, buggy or malicious programs could corrupt other programs
    \item \textbf{Network stack}: Raw network access requires privilege and coordination
    \item \textbf{File system}: File permissions and access control need kernel enforcement
\end{itemize}

The kernel provides \textbf{controlled access} to these resources through system calls.

\subsection{User Mode vs. Kernel Mode}

CPUs support different privilege levels (called \textbf{rings} on x86):
\begin{itemize}
    \item \textbf{Ring 0 (Kernel Mode)}: Unrestricted access to all hardware and memory
    \item \textbf{Ring 3 (User Mode)}: Restricted access; certain instructions cause traps
\end{itemize}

System calls provide a \textbf{controlled, secure gateway} from user mode to kernel mode.

\subsection{How System Calls Work}

Making a system call involves:

\begin{enumerate}
    \item \textbf{Program preparation}: Place syscall number and arguments in registers
    \item \textbf{Trap instruction}: Execute special instruction (e.g., \texttt{syscall} on x86-64)
    \item \textbf{CPU mode switch}: CPU switches from user mode to kernel mode
    \item \textbf{Kernel execution}: Kernel validates arguments and performs the operation
    \item \textbf{Result return}: Kernel places result in a register and switches back to user mode
\end{enumerate}

\textbf{Example: The \texttt{write} syscall on x86-64}

\begin{lstlisting}[language=nasm]
; User-space code making a raw syscall
mov rax, 1          ; Syscall number for 'write' (on x86-64 Linux)
mov rdi, 1          ; File descriptor 1 (stdout)
mov rsi, msg        ; Pointer to message
mov rdx, 13         ; Length of message
syscall             ; Invoke kernel

; Result is now in rax
\end{lstlisting}

System calls are the \textbf{only supported interface} between user and kernel code.

\section{The C Standard Library (libc)}

\subsection{What is libc?}

\textbf{libc} is the C standard library—a collection of functions that implement the C standard library (ISO C) and POSIX APIs. On Linux, the most common implementation is \textbf{glibc} (GNU C Library).

Key responsibilities:
\begin{itemize}
    \item \textbf{Standard C functions}: \texttt{printf}, \texttt{malloc}, \texttt{strcpy}, etc.
    \item \textbf{System call wrappers}: \texttt{open}, \texttt{read}, \texttt{write}, etc.
    \item \textbf{Buffering and caching}: Optimize performance by batching operations
    \item \textbf{Thread safety}: Support for multithreaded programs
    \item \textbf{Locale support}: Internationalization and localization
\end{itemize}

\subsection{libc as a Wrapper Layer}

Most libc functions are \textbf{wrappers} around system calls, adding value through buffering, error handling, and abstraction.

\textbf{Example 1: \texttt{printf}}

\texttt{printf} is \textbf{not} a system call. It's a complex user-space function:

\begin{lstlisting}[language=C]
// Simplified view of what printf does
int printf(const char *format, ...) {
    va_list args;
    va_start(args, format);

    // 1. Format the string in user space
    char buffer[4096];
    int len = vsnprintf(buffer, sizeof(buffer), format, args);

    // 2. Write to stdout (which is buffered!)
    //    stdout is a FILE* with its own buffer
    return fwrite(buffer, 1, len, stdout);

    // fwrite will eventually call the write() syscall,
    // possibly after buffering many small writes
}
\end{lstlisting}

The chain looks like:
\begin{verbatim}
printf()
    → vfprintf() [format the string]
        → fwrite() [manage stdio buffer]
            → write() syscall [actual kernel I/O]
\end{verbatim}

\textbf{Why this indirection?}
\begin{enumerate}
    \item \textbf{Performance}: Formatting in user space avoids kernel mode switches
    \item \textbf{Buffering}: Multiple \texttt{printf} calls can be combined into one \texttt{write} syscall
    \item \textbf{Flexibility}: The same formatting code works for files, strings, and custom streams
\end{enumerate}

\textbf{Example 2: \texttt{malloc}}

\texttt{malloc} is also \textbf{not} a system call (mostly):

\begin{lstlisting}[language=C]
void *malloc(size_t size) {
    // Ask the heap allocator (brk/mmap syscalls under the hood)
    // The heap allocator manages a large memory region
    // and carves it into small pieces for your program

    // Most malloc() calls don't trigger syscalls!
    // They just carve out space from already-allocated regions
    return heap_alloc(size);
}
\end{lstlisting}

The heap allocator (\texttt{malloc} implementation) manages large memory regions obtained from the kernel via \texttt{brk} or \texttt{mmap} syscalls. Most \texttt{malloc} calls just carve out space from these regions—no syscall needed!

\textbf{Example 3: \texttt{strlen}}

\texttt{strlen} is purely user-space:

\begin{lstlisting}[language=C]
size_t strlen(const char *s) {
    size_t len = 0;
    while (s[len]) len++;  // Just reads memory!
    return len;
}
\end{lstlisting}

No syscall involved—it's just reading your program's own memory.

\subsection{When libc Does Make Syscalls}

Some libc functions are thin wrappers around syscalls:

\begin{lstlisting}[language=C]
// glibc's open() function (simplified)
int open(const char *pathname, int flags, mode_t mode) {
    // This is basically just a syscall wrapper
    return syscall(__NR_open, pathname, flags, mode);
}
\end{lstlisting}

But even "simple" wrappers add value:
\begin{itemize}
    \item \textbf{Error handling}: Convert syscall return values to \texttt{errno}
    \item \textbf{Thread safety}: Handle cancellation points for pthreads
    \item \textbf{Cancellation checking}: Allow threads to be cancelled at safe points
\end{itemize}

\section{System Call Numbers and Conventions}

\subsection{Syscall Numbers}

Each system call has a unique number. On Linux x86-64:

\begin{verbatim}
0   read
1   write
2   open
3   close
... hundreds more ...
\end{verbatim}

These numbers are architecture-specific! ARM64 has different numbers than x86-64.

\textbf{Raw syscall example}:

\begin{lstlisting}[language=C]
// Make a write syscall directly (bypassing libc)
#include <unistd.h>
#include <sys/syscall.h>

int main(void) {
    const char msg[] = "Hello via raw syscall!\n";

    // Direct syscall (x86-64 Linux)
    syscall(__NR_write, 1, msg, sizeof(msg) - 1);

    return 0;
}
\end{lstlisting}

\subsection{Syscall Calling Conventions}

On x86-64 Linux:
\begin{itemize}
    \item \textbf{Syscall number}: RAX register
    \item \textbf{Arguments}: RDI, RSI, RDX, R10, R8, R9 (in order)
    \item \textbf{Return value}: RAX register
    \item \textbf{Error indication}: Return value between -4095 and -1 (negated errno)
\end{itemize}

Note the difference from function calls:
\begin{itemize}
    \item Function calls use RCX for 4th argument
    \item Syscalls use R10 for 4th argument (syscall instruction clobbers RCX)
\end{itemize}

\section{Error Reporting: errno}

\subsection{How errno Works}

When a syscall fails, the kernel returns a negative error code. libc converts this to:

\begin{enumerate}
    \item \textbf{Return value}: -1 (for most functions)
    \item \textbf{errno}: A global variable containing the error code
\end{enumerate}

\begin{lstlisting}[language=C]
#include <stdio.h>
#include <errno.h>
#include <string.h>

int main(void) {
    FILE *f = fopen("/nonexistent/file", "r");

    if (f == NULL) {
        // errno is set by libc based on kernel error
        printf("Error %d: %s\n", errno, strerror(errno));
    }

    return 0;
}
\end{lstlisting}

\textbf{Important}: \texttt{errno} is thread-local! Each thread has its own copy, so multithreaded programs work correctly.

\subsection{Why errno Exists}

Why not just return the error code directly?

\begin{enumerate}
    \item \textbf{Historical}: Early C had limited error handling options
    \item \textbf{Efficiency}: Negative return values can indicate valid data (e.g., \texttt{read} can legitimately return 0 for EOF)
    \item \textbf{Detail}: errno provides more error information than just "success/failure"
\end{enumerate}

\section{Observing Syscalls with strace}

\texttt{strace} is a powerful tool that traces system calls made by a program.

\subsection{Basic strace Usage}

Create \texttt{simple.c}:
\begin{lstlisting}[language=C]
#include <stdio.h>
#include <stdlib.h>
#include <string.h>

int main(void) {
    char *s = strdup("Hello");
    printf("%s\n", s);
    free(s);
    return 0;
}
\end{lstlisting}

Compile and trace:
\begin{lstlisting}[language=bash]
gcc -o simple simple.c
strace ./simple
\end{lstlisting}

Typical output (abbreviated):
\begin{verbatim}
execve("./simple", ["./simple"], 0x7ffc...) = 0
brk(NULL)                               = 0x55555000
brk(0x55556c000)                        = 0x55556c000
mmap(NULL, 4096, PROT_READ|PROT_WRITE, MAP_PRIVATE|MAP_ANONYMOUS, -1, 0) = 0x7ffff...
write(1, "Hello\n", 6)                  = 6
munmap(0x7ffff..., 4096)                = 0
exit_group(0)                           = ?
+++ exited with 0 +++
\end{verbatim}

\textbf{What to observe}:
\begin{itemize}
    \item \texttt{brk}/\texttt{mmap}: Memory allocation for heap
    \item \texttt{write}: The actual I/O syscall (only one, despite printf!)
    \item No syscall for \texttt{strlen}-like operations (they're pure user-space)
\end{itemize}

\section{Library vs. Syscall Performance}

\subsection{The Cost of Syscalls}

System calls are expensive compared to function calls:

\begin{verbatim}
Function call: ~1-5 nanoseconds
System call: ~500-2000 nanoseconds
\end{verbatim}

The overhead comes from:
\begin{enumerate}
    \item \textbf{Mode switch}: User → kernel → user transition
    \item \textbf{Validation}: Kernel must validate all arguments
    \item \textbf{Scheduling}: Other processes might run while in kernel
    \item \textbf{Cache effects}: Kernel code uses different cache lines
\end{enumerate}

\subsection{libc Buffering}

libc performs buffering at several levels:

\textbf{Stdio buffering (\texttt{FILE*} streams)}:
\begin{lstlisting}[language=C]
// Unbuffered (slow)
setbuf(stdout, NULL);
for (int i = 0; i < 1000; i++) {
    printf("x");  // Syscall each iteration!
}

// Fully buffered (fast)
setvbuf(stdout, NULL, _IOFBF, BUFSIZ);
for (int i = 0; i < 1000; i++) {
    printf("x");  // Buffered until:
                 // - Buffer is full
                 // - fflush() is called
                 // - Program exits
}
\end{lstlisting}

\textbf{Default buffering behavior}:
\begin{itemize}
    \item \textbf{stdin}: Line-buffered (when interactive) or fully buffered
    \item \textbf{stdout}: Line-buffered (when connected to terminal) or fully buffered
    \item \textbf{stderr}: Always unbuffered (ensure error messages appear)
\end{itemize}

\section{Key Takeaways}

\begin{enumerate}
    \item \textbf{libc vs. kernel}: Most C library functions are user-space code that may eventually trigger syscalls, but many operations (like strlen, malloc from existing heap, etc.) don't require kernel interaction.
    \item \textbf{Syscalls are expensive}: System calls cost $\sim$500-2000 nanoseconds, much more than function calls. This is why libc buffering is so important for performance.
    \item \textbf{Buffering matters}: libc's stdio buffering can reduce syscall count by orders of magnitude, dramatically improving performance for I/O operations.
    \item \textbf{strace is powerful}: The strace tool reveals exactly which syscalls a program makes, invaluable for debugging and performance analysis.
    \item \textbf{Error reporting}: The kernel returns error codes in registers; libc converts these to errno for easier error handling.
    \item \textbf{Memory management}: malloc/free mostly manage user-space heap, only calling brk/mmap when more memory is needed from the kernel.
    \item \textbf{User/kernel boundary}: Understanding this boundary is essential for writing efficient, portable C programs and for debugging system-level issues.
\end{enumerate}

\section{Looking Ahead}

In Chapter 9, we'll explore \textbf{ABI and cross-platform compilation}—how different architectures and operating systems define binary interfaces, and what it takes to make code portable across different platforms.
